\section{优化与 GRATE：从大模到精简模型}
\subsection{两 Agent + 质量指标的训练流水线}
在 00:24:31--00:31:15 的案例中，ServiceNow 先用参数量达 405B 的 Llama 3.5 大模型打造高质量交互，再将 tape 记录用作数据，微调参数仅 8B 的 Llama 模型以恢复性能。结果是：微调模型的 GRATE 分数与原始大模型等价，却在运行成本上实现 300x+ 的节省。

\begin{knowledgebox}{GRATE 质量得分}
GRATE 是一套可量化的质量属性（grounded、responsive、accurate、transparent、effective），每一维度都可独立测量与优化，方便将 tape 作为训练时的 reward/regression signal。
\end{knowledgebox}

\subsection{成本-性能权衡}
指标图在 00:30:25--00:31:38 显示：x 轴代表每百万轮对话的成本，y 轴代表 GRATE 得分；大型模型仍在右上角，但通过 tape 生成的微调模型可以跑出接近的质量却只需千分之一的费用。

\subsection{本章小结}
TapeAgents 不仅便于工程化开发，也让我们可以把昂贵的“老师级” Agent 转化为高质量且低成本的“学生 Agent”，成本/效果曲线可以通过 GRATE 可视化。

